% titel: Blatt 7, A6: Ideale Gasgleichung in vier Varianten
% skript: Mortimer Kap. 11 (ideales Gasgesetz)

% Teil I: Werkzeuge, kein L\"osungsweg.

Das ideale Gasgesetz $pV = nRT$ verbindet Druck, Volumen, Stoffmenge und
Temperatur. Je nach c) wird daraus die Konzentration $c=n/V=p/(RT)$.

\begin{werkzeugkasten}[Ideal-Gas-Formeln]
\begin{itemize}
  \item $pV = nRT$, $R=0.0821\ \mathrm{L\,atm\,mol^{-1}K^{-1}}$
        (oder $8.314\ \mathrm{J\,mol^{-1}K^{-1}}$, passend zu Pa/Pa$\cdot$m$^3$).
  \item Konzentration: $c = \frac{n}{V} = \frac{p}{RT}$.
  \item Molvolumen: $V_m = \frac{V}{n}=\frac{RT}{p}$.
  \item Bei konstantem Volumen: $p\propto T$; bei konstanter Temperatur:
        $pV = \text{const}$ (Boyle).
\end{itemize}
\textit{Skript:} Mortimer Kap. 11, ideales Gasgesetz.
\end{werkzeugkasten}

\begin{achtung}[Darauf achten]
\begin{itemize}
  \item Kelvin: $0\,^\circ\mathrm{C} = 273.15\,\mathrm{K}$.
  \item Einheiten koppeln: Mit $R$ in J/(mol$\cdot$K) und Druck in Pa
        folgt $c$ in $\mathrm{mol\,m^{-3}}$ (nicht L$^{-1}$).
  \item Standardbedingungen: $0\,^\circ\mathrm{C}$ und $1$ atm
        $\Rightarrow V_m \approx 22.4\ \mathrm{L\,mol^{-1}}$.
\end{itemize}
\end{achtung}
